\documentclass[11pt]{article}

\usepackage[a4paper,margin=1in]{geometry}
\usepackage{amsmath,amssymb,amsthm,mathtools}
\usepackage{enumitem}
\usepackage{hyperref}

\newtheorem{theorem}{Theorem}[section]
\newtheorem{lemma}[theorem]{Lemma}
\newtheorem{corollary}[theorem]{Corollary}
\newtheorem{definition}[theorem]{Definition}
\newtheorem{remark}[theorem]{Remark}

\DeclareMathOperator{\cl}{cl}
\DeclareMathOperator{\spec}{Spec}
\DeclareMathOperator{\rad}{rad}
\DeclareMathOperator{\tr}{tr}

\newcommand{\R}{\mathbb{R}}
\newcommand{\N}{\mathbb{N}}
\newcommand{\K}{\mathcal{K}}
\newcommand{\V}{V}
\newcommand{\E}{E}
\newcommand{\bbeta}{\beta}
\newcommand{\ol}{\overline}

\title{Fisher's lower bound for triangle density on \(1/2\le p\le 2/3\)}
\author{}
\date{}

\begin{document}

\maketitle

\begin{abstract}
We give a detailed proof of Fisher's lower bound for the number of
triangles in a finite graph with edge density between \(1/2\) and \(2/3\).
The proof uses the dependence polynomial of a graph, its smallest positive
zero, a spectral lower bound for the corresponding growth factor, and an
elementary one-variable optimisation.  The presentation is organised so that
each nontrivial step can be translated into a separate Lean lemma.
\end{abstract}

\section{Statement of the result}

\begin{theorem}[Fisher's triangle lower bound]
\label{thm:fisher-finite}
Let \(G\) be a finite simple graph with
\[
        n := |\V(G)|,\qquad e := |\E(G)|,\qquad T := \#\{K_3\subseteq G\}.
\]
Assume
\[
        \frac{n^2}{4}\le e\le \frac{n^2}{3}.
\]
Then
\[
        T
        \ge
        \frac{9en-2n^3-2\bigl(n^2-3e\bigr)^{3/2}}{27}.
\]
\end{theorem}

\begin{corollary}[Density form]
\label{cor:density-form}
Let
\[
        p := \frac{2e}{n^2},
        \qquad
        q := \frac{6T}{n^3}.
\]
If \(1/2\le p\le 2/3\), then
\[
        q
        \ge
        p-\frac49-\frac49\left(1-\frac{3p}{2}\right)^{3/2}.
\]
\end{corollary}

\section{Graph-theoretic and polynomial preliminaries}

\begin{definition}[Clique counts]
For \(k\ge 0\), let \(c_k(G)\) be the number of \(k\)-cliques of \(G\).
We include the empty clique, so \(c_0(G)=1\).  In particular,
\[
        c_1(G)=n,\qquad c_2(G)=e,\qquad c_3(G)=T.
\]
\end{definition}

\begin{definition}[Dependence polynomial]
The dependence polynomial of \(G\) is
\[
        D_G(z)
        :=
        \sum_{k\ge 0}(-1)^k c_k(G)z^k.
\]
Equivalently,
\[
        D_G(z)=1-nz+ez^2-Tz^3+c_4(G)z^4-\cdots.
\]
\end{definition}

\begin{definition}[Common neighbourhood of a clique]
If \(S\subseteq \V(G)\) is a clique, define
\[
        N_G(S)
        :=
        \{v\in \V(G)\setminus S:\ v\text{ is adjacent to every element of }S\}.
\]
Let \(G[N_G(S)]\) be the induced subgraph on this vertex set.
\end{definition}

\begin{lemma}[Derivative identity for the dependence polynomial]
\label{lem:derivative-identity}
Let \(j\ge 0\).  Then
\[
        D_G^{(j)}(z)
        =
        (-1)^j j!
        \sum_{\substack{S\subseteq \V(G)\\ S\text{ is a }j\text{-clique}}}
        D_{G[N_G(S)]}(z).
\]
For \(j=0\), this is interpreted as the identity \(D_G=D_G\).
\end{lemma}

\begin{proof}
For \(j\ge 1\),
\[
        D_G(z)=\sum_{C\in \K(G)}(-z)^{|C|},
\]
where \(\K(G)\) is the set of all cliques of \(G\), including the empty clique.
Differentiating \(j\) times gives
\[
        D_G^{(j)}(z)
        =
        \sum_{C\in \K(G),\, |C|\ge j}
        (-1)^{|C|}
        |C|(|C|-1)\cdots(|C|-j+1)
        z^{|C|-j}.
\]
Since
\[
        |C|(|C|-1)\cdots(|C|-j+1)
        =
        j!\binom{|C|}{j},
\]
we may rewrite this as
\[
        D_G^{(j)}(z)
        =
        (-1)^j j!
        \sum_{C\in \K(G),\, |C|\ge j}
        \binom{|C|}{j}
        (-z)^{|C|-j}.
\]
Counting a pair \((S,C)\), where \(S\subseteq C\), \(|S|=j\), and \(C\) is a clique,
instead by first choosing \(S\), gives
\[
        D_G^{(j)}(z)
        =
        (-1)^j j!
        \sum_{\substack{S\in \K(G)\\ |S|=j}}
        \sum_{\substack{R\in \K(G[N_G(S)])}}
        (-z)^{|R|}.
\]
The inner sum is precisely \(D_{G[N_G(S)]}(z)\).
\end{proof}

\section{The root and growth lemmas}

The following facts are the algebraic core of Fisher's method.

\begin{definition}[Graph monoid]
Let \(M(G)\) be the monoid generated by the vertices of \(G\), with relations
\[
        uv=vu
        \qquad\text{whenever } uv\in \E(G).
\]
Let \(a_k(G)\) be the number of equivalence classes in \(M(G)\) represented by
words of length \(k\).
\end{definition}

\begin{lemma}[Cartier--Foata inversion formula]
\label{lem:cartier-foata}
As a formal power series,
\[
        \sum_{k\ge 0}a_k(G)z^k
        =
        \frac{1}{D_G(z)}.
\]
\end{lemma}

\begin{proof}
This is the standard Cartier--Foata inversion for trace monoids.  We recall the
argument because it is useful for formalisation.

Let a trace be an equivalence class of words in the graph monoid \(M(G)\).  A clique
\(C\subseteq \V(G)\) may be viewed as a commutative product of its vertices.  Consider
pairs \((m,C)\), where \(m\in M(G)\) is a trace and \(C\) is a clique.  Give the pair
weight
\[
        \mathrm{wt}(m,C):=(-1)^{|C|}z^{|m|+|C|}.
\]
The coefficient of \(z^r\) in
\[
        \left(\sum_{k\ge 0}a_k(G)z^k\right)D_G(z)
\]
is the signed weight of all pairs \((m,C)\) with \(|m|+|C|=r\).

For \(r>0\), define a sign-reversing involution as follows.  Put a fixed total order
on \(\V(G)\).  Among all vertices that can be moved to the right end of a representative
of \(m\), together with the vertices of \(C\), choose the largest one.  If this vertex lies
in \(C\), move it from \(C\) to the right end of \(m\).  If it lies at the right end of \(m\),
move it from \(m\) into \(C\).  The clique condition is preserved because exactly the
vertices movable to the right end commute pairwise.  The operation changes \(|C|\) by
one and preserves total degree, so the two paired terms cancel.

The only unpaired object occurs in degree \(0\), namely \((1,\varnothing)\).  Hence the
product is \(1\), proving the formal identity.
\end{proof}

\begin{lemma}[Existence of the smallest positive zero]
\label{lem:smallest-positive-zero}
There exists a number \(\beta_G>0\) such that
\[
        D_G(\beta_G)=0,
\]
and \(D_G(x)>0\) for all \(0\le x<\beta_G\).  Moreover, \(\beta_G\) is the radius of
convergence of \(1/D_G(z)\).
\end{lemma}

\begin{proof}
By Lemma~\ref{lem:cartier-foata},
\[
        \frac{1}{D_G(z)}=\sum_{k\ge 0}a_k(G)z^k
\]
has nonnegative coefficients.  Since \(G\) is finite, the sequence \(a_k(G)\) has at
most exponential growth and at least one nonzero term in every degree, so the radius
of convergence is a positive real number.  By Pringsheim's theorem for power series
with nonnegative coefficients, the radius of convergence is a singularity on the positive
real axis.  Since \(D_G\) is a polynomial and \(1/D_G\) is rational, this singularity is a
zero of \(D_G\).  Denote it by \(\beta_G\).

Because \(D_G(0)=1\), continuity implies that \(D_G\) remains positive on the interval
from \(0\) up to its first positive zero.  Therefore \(D_G(x)>0\) for \(0\le x<\beta_G\).
\end{proof}

\begin{definition}[Growth factor]
Define
\[
        r(G):=\frac{1}{\beta_G}.
\]
By the Cauchy--Hadamard formula and Lemma~\ref{lem:cartier-foata},
\[
        r(G)=\limsup_{k\to\infty} a_k(G)^{1/k}.
\]
\end{definition}

\begin{lemma}[Induced-subgraph monotonicity]
\label{lem:induced-monotonicity}
If \(H\) is an induced subgraph of \(G\), then
\[
        \beta_G\le \beta_H.
\]
Consequently,
\[
        D_H(x)\ge 0
        \qquad\text{for every }x\in[0,\beta_G].
\]
\end{lemma}

\begin{proof}
The inclusion \(\V(H)\subseteq \V(G)\) induces an injective map from traces in \(M(H)\)
to traces in \(M(G)\), because the commutation relations among vertices of \(H\) are
exactly the same as their commutation relations in \(G\).  Therefore
\[
        a_k(H)\le a_k(G)
        \qquad\text{for every }k.
\]
Taking limsups gives
\[
        r(H)
        =
        \limsup_{k\to\infty}a_k(H)^{1/k}
        \le
        \limsup_{k\to\infty}a_k(G)^{1/k}
        =
        r(G).
\]
Since \(r(X)=1/\beta_X\), this is equivalent to \(\beta_G\le \beta_H\).

Now let \(x\in[0,\beta_G]\).  Since \(\beta_G\le\beta_H\), we have \(x\le\beta_H\).
By Lemma~\ref{lem:smallest-positive-zero}, \(D_H\) is positive on \([0,\beta_H)\) and
vanishes at \(\beta_H\).  Hence \(D_H(x)\ge 0\).
\end{proof}

\section{Fisher's cubic-root inequality}

\begin{lemma}[Third-order truncation inequality]
\label{lem:third-truncation}
Let
\[
        D_G(z)=1-nz+ez^2-Tz^3+c_4z^4-\cdots.
\]
Then
\[
        1-n\beta_G+e\beta_G^2-T\beta_G^3\le 0.
\]
Equivalently,
\[
        r(G)^3-nr(G)^2+er(G)-T\le 0.
\]
\end{lemma}

\begin{proof}
Apply Taylor's theorem with integral remainder at order \(3\):
\[
        D_G(\beta_G)
        =
        D_G(0)+D_G'(0)\beta_G+\frac{D_G''(0)}{2}\beta_G^2
        +\frac{D_G^{(3)}(0)}{6}\beta_G^3
        +
        \frac{1}{6}\int_0^{\beta_G}(\beta_G-t)^3D_G^{(4)}(t)\,dt.
\]
The first four Taylor terms are
\[
        1-n\beta_G+e\beta_G^2-T\beta_G^3.
\]
Since \(D_G(\beta_G)=0\), it remains to show that the integral remainder is nonnegative.

By Lemma~\ref{lem:derivative-identity} with \(j=4\),
\[
        D_G^{(4)}(t)
        =
        4!
        \sum_{\substack{S\subseteq\V(G)\\S\text{ is a }4\text{-clique}}}
        D_{G[N_G(S)]}(t).
\]
For each \(4\)-clique \(S\), the graph \(G[N_G(S)]\) is an induced subgraph of \(G\).
Therefore, by Lemma~\ref{lem:induced-monotonicity},
\[
        D_{G[N_G(S)]}(t)\ge 0
        \qquad\text{for all }t\in[0,\beta_G].
\]
Hence \(D_G^{(4)}(t)\ge 0\) on \([0,\beta_G]\).  The kernel
\((\beta_G-t)^3/6\) is also nonnegative on this interval, so the integral remainder is
nonnegative.  Therefore
\[
        0
        =
        D_G(\beta_G)
        \ge
        1-n\beta_G+e\beta_G^2-T\beta_G^3.
\]
Finally, multiplying by \(r(G)^3=\beta_G^{-3}\) gives
\[
        r(G)^3-nr(G)^2+er(G)-T\le 0.
\]
\end{proof}

\begin{lemma}[Cubic root consequence]
\label{lem:cubic-root-consequence}
Let
\[
        \phi(x):=x^3-nx^2+ex.
\]
Then there exists a real number \(R\ge r(G)\) such that
\[
        T=\phi(R).
\]
\end{lemma}

\begin{proof}
Define
\[
        P(x):=x^3-nx^2+ex-T.
\]
By Lemma~\ref{lem:third-truncation},
\[
        P(r(G))\le 0.
\]
Since \(P(x)\to +\infty\) as \(x\to+\infty\), the intermediate value theorem gives
a real root \(R\in[r(G),\infty)\).  For this root,
\[
        0=P(R)=R^3-nR^2+eR-T,
\]
so
\[
        T=R^3-nR^2+eR=\phi(R).
\]
\end{proof}

\section{Fisher's spectral lower bound on \(r(G)\)}

\begin{definition}[Complement spectral radius]
Let \(A_{\overline G}\) be the adjacency matrix of the complement graph \(\overline G\).
Let
\[
        \lambda(\overline G)
\]
denote the largest eigenvalue of \(A_{\overline G}\).  Since \(A_{\overline G}\) is a
real symmetric nonnegative matrix, this is its spectral radius.
\end{definition}

\begin{lemma}[Spectral growth bound]
\label{lem:spectral-growth-bound}
For every finite simple graph \(G\),
\[
        r(G)\ge 1+\lambda(\overline G).
\]
\end{lemma}

\begin{proof}
Let \(B:=I+A_{\overline G}\).  A word
\[
        v_1v_2\cdots v_k
\]
is called reduced-adjacent if, for every \(i\),
\[
        v_i=v_{i+1}
        \quad\text{or}\quad
        v_iv_{i+1}\notin \E(G).
\]
Equivalently, \((v_i,v_{i+1})\) is an edge of the looped complement graph whose
adjacency matrix is \(B\).

Such a word admits no adjacent commutation move in the graph monoid \(M(G)\): two
distinct adjacent letters commute exactly when they are adjacent in \(G\).  Therefore each
reduced-adjacent word represents a singleton trace class.  Hence the number of traces
of length \(k\) is at least the number of length-\((k-1)\) walks in the looped complement:
\[
        a_k(G)
        \ge
        \mathbf{1}^{\mathsf T}B^{k-1}\mathbf{1}
        \qquad(k\ge 1).
\]
Taking exponential growth rates gives
\[
        r(G)
        =
        \limsup_{k\to\infty}a_k(G)^{1/k}
        \ge
        \rho(B),
\]
where \(\rho(B)\) is the spectral radius of \(B\).  Since \(B=I+A_{\overline G}\),
its eigenvalues are \(1+\mu\), where \(\mu\) ranges over the eigenvalues of
\(A_{\overline G}\).  Thus
\[
        \rho(B)=1+\lambda(\overline G).
\]
\end{proof}

\begin{lemma}[Average-degree spectral bound]
\label{lem:average-degree-spectral}
Let \(m:=|\E(\overline G)|\).  Then
\[
        \lambda(\overline G)\ge \frac{2m}{n}.
\]
Consequently,
\[
        r(G)\ge n-\frac{2e}{n}.
\]
\end{lemma}

\begin{proof}
By the Rayleigh quotient applied to the all-ones vector \(\mathbf{1}\),
\[
        \lambda(\overline G)
        \ge
        \frac{\mathbf{1}^{\mathsf T}A_{\overline G}\mathbf{1}}
             {\mathbf{1}^{\mathsf T}\mathbf{1}}
        =
        \frac{2m}{n}.
\]
Since
\[
        m=\binom n2-e=\frac{n(n-1)}2-e,
\]
Lemma~\ref{lem:spectral-growth-bound} gives
\[
        r(G)
        \ge
        1+\frac{2m}{n}
        =
        1+\frac{n(n-1)-2e}{n}
        =
        n-\frac{2e}{n}.
\]
\end{proof}

\section{The one-variable optimisation}

\begin{lemma}[Critical points of the cubic]
\label{lem:cubic-critical-points}
Let
\[
        \phi(x):=x^3-nx^2+ex,
        \qquad
        s:=\sqrt{n^2-3e}.
\]
If \(e\le n^2/3\), then \(s\) is real and the critical points of \(\phi\) are
\[
        x_-:=\frac{n-s}{3},
        \qquad
        x_+:=\frac{n+s}{3}.
\]
Moreover,
\[
        \phi'(x)\ge 0 \text{ on }(-\infty,x_-]\cup[x_+,\infty),
        \qquad
        \phi'(x)\le 0 \text{ on }[x_-,x_+].
\]
Therefore \(\phi\) attains its minimum on \([x_-,\infty)\) at \(x_+\).
\end{lemma}

\begin{proof}
We have
\[
        \phi'(x)=3x^2-2nx+e.
\]
The discriminant of this quadratic is
\[
        (-2n)^2-12e=4(n^2-3e)=4s^2.
\]
Hence the two roots are
\[
        \frac{2n-2s}{6}=\frac{n-s}{3}=x_-,
        \qquad
        \frac{2n+2s}{6}=\frac{n+s}{3}=x_+.
\]
Since the leading coefficient \(3\) is positive, \(\phi'\) is nonnegative outside the
interval \([x_-,x_+]\) and nonpositive inside it.  Hence \(\phi\) increases on
\((-\infty,x_-]\), decreases on \([x_-,x_+]\), and increases on \([x_+,\infty)\).
Thus its minimum on \([x_-,\infty)\) is attained at \(x_+\).
\end{proof}

\begin{lemma}[The spectral lower bound lies beyond \(x_-\)]
\label{lem:spectral-beyond-xminus}
Assume \(e\le n^2/3\), and let
\[
        s:=\sqrt{n^2-3e},
        \qquad
        x_-:=\frac{n-s}{3}.
\]
Then
\[
        n-\frac{2e}{n}\ge x_-.
\]
\end{lemma}

\begin{proof}
Since \(s^2=n^2-3e\), we have
\[
        e=\frac{n^2-s^2}{3}.
\]
Therefore
\[
        n-\frac{2e}{n}
        =
        n-\frac{2(n^2-s^2)}{3n}
        =
        \frac{n}{3}+\frac{2s^2}{3n}.
\]
Thus
\[
        n-\frac{2e}{n}-x_-
        =
        \left(\frac{n}{3}+\frac{2s^2}{3n}\right)
        -
        \frac{n-s}{3}
        =
        \frac{s}{3}+\frac{2s^2}{3n}
        \ge 0.
\]
\end{proof}

\begin{lemma}[Value of the cubic at the lower critical point]
\label{lem:cubic-value}
Let
\[
        s:=\sqrt{n^2-3e},
        \qquad
        x_+:=\frac{n+s}{3}.
\]
Then
\[
        \phi(x_+)
        =
        \frac{9en-2n^3-2(n^2-3e)^{3/2}}{27}.
\]
\end{lemma}

\begin{proof}
Since \(x_+\) is a critical point of \(\phi\), we have
\[
        3x_+^2-2nx_+ + e=0,
\]
so
\[
        e=2nx_+-3x_+^2.
\]
Therefore
\[
        \phi(x_+)
        =
        x_+^3-nx_+^2+ex_+
        =
        x_+^3-nx_+^2+(2nx_+-3x_+^2)x_+
        =
        nx_+^2-2x_+^3.
\]
Hence
\[
        \phi(x_+)=x_+^2(n-2x_+).
\]
Substituting \(x_+=(n+s)/3\), we get
\[
        \phi(x_+)
        =
        \frac{(n+s)^2}{9}\cdot \frac{n-2s}{3}
        =
        \frac{(n+s)^2(n-2s)}{27}.
\]
Expanding the numerator,
\[
        (n+s)^2(n-2s)
        =
        n^3-3ns^2-2s^3.
\]
Since \(s^2=n^2-3e\),
\[
        n^3-3ns^2-2s^3
        =
        n^3-3n(n^2-3e)-2s^3
        =
        9en-2n^3-2s^3.
\]
Finally,
\[
        s^3=(n^2-3e)^{3/2},
\]
which proves the claimed formula.
\end{proof}

\section{Proof of Fisher's finite theorem}

\begin{proof}[Proof of Theorem~\ref{thm:fisher-finite}]
Let
\[
        \phi(x):=x^3-nx^2+ex.
\]
By Lemma~\ref{lem:cubic-root-consequence}, there exists \(R\ge r(G)\) such that
\[
        T=\phi(R).
\]
By Lemma~\ref{lem:average-degree-spectral},
\[
        r(G)\ge n-\frac{2e}{n}.
\]
By Lemma~\ref{lem:spectral-beyond-xminus},
\[
        n-\frac{2e}{n}\ge x_-,
        \qquad
        x_-:=\frac{n-\sqrt{n^2-3e}}{3}.
\]
Therefore
\[
        R\ge x_-.
\]
By Lemma~\ref{lem:cubic-critical-points}, the function \(\phi\) attains its minimum
on \([x_-,\infty)\) at
\[
        x_+:=\frac{n+\sqrt{n^2-3e}}{3}.
\]
Hence
\[
        T=\phi(R)\ge \phi(x_+).
\]
By Lemma~\ref{lem:cubic-value},
\[
        \phi(x_+)
        =
        \frac{9en-2n^3-2(n^2-3e)^{3/2}}{27}.
\]
Thus
\[
        T
        \ge
        \frac{9en-2n^3-2(n^2-3e)^{3/2}}{27}.
\]
\end{proof}

\section{Density normalisation}

\begin{proof}[Proof of Corollary~\ref{cor:density-form}]
Let
\[
        p:=\frac{2e}{n^2},
        \qquad
        q:=\frac{6T}{n^3}.
\]
The assumption \(1/2\le p\le 2/3\) is equivalent to
\[
        \frac{n^2}{4}\le e\le \frac{n^2}{3}.
\]
By Theorem~\ref{thm:fisher-finite},
\[
        T
        \ge
        \frac{9en-2n^3-2(n^2-3e)^{3/2}}{27}.
\]
Multiplying by \(6/n^3\), we obtain
\[
        q
        \ge
        \frac{6}{27}
        \left(
            \frac{9e}{n^2}
            -2
            -2\left(1-\frac{3e}{n^2}\right)^{3/2}
        \right).
\]
Since \(e/n^2=p/2\),
\[
        \frac{9e}{n^2}=\frac{9p}{2},
        \qquad
        1-\frac{3e}{n^2}=1-\frac{3p}{2}.
\]
Therefore
\[
        q
        \ge
        \frac{2}{9}
        \left(
            \frac{9p}{2}
            -2
            -2\left(1-\frac{3p}{2}\right)^{3/2}
        \right).
\]
Simplifying gives
\[
        q
        \ge
        p-\frac49-\frac49\left(1-\frac{3p}{2}\right)^{3/2}.
\]
\end{proof}

\section{Sharpness}

\begin{remark}[Extremal construction]
The bound is asymptotically sharp.  Let
\[
        u:=\sqrt{1-\frac{3p}{2}},
        \qquad 0\le u\le \frac12.
\]
Consider the complete tripartite graphon with part weights
\[
        a=a':=\frac{1+u}{3},
        \qquad
        b:=\frac{1-2u}{3}.
\]
Its edge density is
\[
        1-a^2-a'^2-b^2
        =
        1-2\left(\frac{1+u}{3}\right)^2
          -\left(\frac{1-2u}{3}\right)^2
        =
        \frac{2}{3}(1-u^2)
        =
        p.
\]
Its triangle density is
\[
        6aa'b
        =
        6\left(\frac{1+u}{3}\right)^2
          \left(\frac{1-2u}{3}\right).
\]
Expanding,
\[
        6aa'b
        =
        \frac{2}{9}(1+u)^2(1-2u)
        =
        \frac{2}{9}(1-3u^2-2u^3).
\]
Since \(u^2=1-\frac{3p}{2}\), this becomes
\[
        p-\frac49-\frac49u^3
        =
        p-\frac49-\frac49\left(1-\frac{3p}{2}\right)^{3/2}.
\]
Thus the lower bound is attained asymptotically by complete tripartite graphs whose
part ratios tend to
\[
        \frac{1+u}{3},\quad
        \frac{1+u}{3},\quad
        \frac{1-2u}{3}.
\]
\end{remark}

\section{Lean-formalisation blueprint}

\begin{remark}[Suggested Lean decomposition]
A Lean development could be organised into the following modules.

\begin{enumerate}[label=\textbf{Module \arabic*.}, leftmargin=2.5cm]

\item \textbf{Finite graphs and clique counts.}
Define finite simple graphs, complements, induced subgraphs, cliques, \(c_k(G)\),
edge count, and triangle count.

\item \textbf{Dependence polynomial.}
Define
\[
        D_G(z)=\sum_k(-1)^k c_k(G)z^k.
\]
Prove the derivative identity
\[
        D_G^{(j)}(z)
        =
        (-1)^j j!
        \sum_{S\in K_j(G)}
        D_{G[N_G(S)]}(z).
\]
This is purely finite combinatorics plus polynomial differentiation.

\item \textbf{Trace monoid growth.}
Define the graph monoid and prove the Cartier--Foata identity
\[
        \sum_k a_k(G)z^k = D_G(z)^{-1}.
\]
For a first Lean version, this can be imported as a theorem and later formalised by the
sign-reversing involution used above.

\item \textbf{Smallest positive root and monotonicity.}
Using nonnegative coefficients and the radius of convergence, define \(\beta_G\), prove
\(D_G(x)>0\) for \(0\le x<\beta_G\), and prove induced-subgraph monotonicity
\[
        H\subseteq_{\mathrm{ind}}G
        \implies
        \beta_G\le\beta_H.
\]

\item \textbf{Third truncation.}
Use Taylor's theorem with integral remainder and the fourth-derivative identity to prove
\[
        1-n\beta_G+e\beta_G^2-T\beta_G^3\le 0.
\]

\item \textbf{Spectral estimate.}
Define the adjacency matrix of \(\overline G\), prove
\[
        r(G)\ge 1+\lambda(\overline G),
\]
by injecting looped-complement walks into trace classes, and prove
\[
        \lambda(\overline G)\ge \frac{2|\E(\overline G)|}{n}
\]
by the Rayleigh quotient.

\item \textbf{Real calculus.}
For
\[
        \phi(x)=x^3-nx^2+ex,
\]
prove the critical-point calculation and the explicit value
\[
        \phi\left(\frac{n+\sqrt{n^2-3e}}{3}\right)
        =
        \frac{9en-2n^3-2(n^2-3e)^{3/2}}{27}.
\]

\item \textbf{Final theorem.}
Combine the cubic-root consequence, spectral estimate, and calculus minimisation.

\end{enumerate}
\end{remark}

\end{document}